\documentclass[10pt,aspectratio=169]{beamer}
\usepackage[utf8]{inputenc}
\usepackage[spanish]{babel}
\usepackage{xcolor}
\usepackage{booktabs}

\usetheme{Madrid}
\usecolortheme{default}

\definecolor{democolor}{RGB}{0,100,0}
\newcommand{\demo}[1]{\textcolor{democolor}{\texttt{#1}}}

\title[El Juego de la Vida]{El Juego de la Vida de Conway}
\subtitle{Programación Paralela y Concurrente -- 2026}
\author[Equipo Rosa]{
  L. Guadalupe Jimenez G. \and
  Lázaro García \and
  Israel Olvera Juárez \and
  Jorge F. Zarate D. \and
  Miguel Ángel Torres P. \and
  Andrés Soria C.
}
\institute[UNAM]{Universidad Nacional Autónoma de México}
\date{Mayo 2026}

\begin{document}

% ================================================================
% 1. PORTADA
% ================================================================
\begin{frame}
  \titlepage
\end{frame}

% ================================================================
% 2. INTRODUCCIÓN -- Lupita
% ================================================================
\begin{frame}{¿Qué es el Juego de la Vida?}

  \vspace{0.4cm}

  \begin{itemize}
    \item Creado por \textbf{John Conway} en 1970.
    \item Un \textbf{autómata celular}: grilla de células que evolucionan
          solas según reglas simples.
    \item Con solo 4 reglas emergen patrones sorprendentemente complejos.
    \item Es \textbf{Turing-completo}: puede hacer cualquier cálculo.
  \end{itemize}

  \vspace{0.5cm}

  \begin{center}
    \textbf{Demo en vivo}

    \demo{blinker} \quad \demo{block}

    \small Ejecutar: \texttt{python juego\_vida\_optimo.py}
  \end{center}
\end{frame}

% ================================================================
% REGLAS DEL JUEGO
% ================================================================
\begin{frame}{Las 4 reglas de Conway}

\begin{center}
\Large Toda la simulación surge de solo 4 reglas
\end{center}

\vspace{0.4cm}

\begin{block}{Reglas}
\begin{enumerate}
    \item Una célula viva con menos de 2 vecinos muere.
    \item Una célula viva con 2 o 3 vecinos sobrevive.
    \item Una célula viva con más de 3 vecinos muere.
    \item Una célula muerta con exactamente 3 vecinos nace.
\end{enumerate}
\end{block}

\vspace{0.4cm}

\begin{alertblock}{Idea importante}
Comportamientos extremadamente complejos emergen
a partir de reglas extremadamente simples.
\end{alertblock}

\end{frame}
% ================================================================
% 4. STILL LIFES Y OSCILADORES -- Lázaro
% ================================================================
\begin{frame}{Patrones que persisten en el tiempo}

  \vspace{0.4cm}

  \begin{block}{Still lifes -- no cambian}
    \begin{itemize}
      \item Son estructuras \textbf{estables}: cada célula tiene 2 o 3 vecinos.
      \item Aparecen naturalmente del caos y funcionan como ``ladrillos''.
    \end{itemize}
  \end{block}

  \begin{block}{Osciladores -- alternan entre estados}
    \begin{itemize}
      \item Ciclan entre 2 o más configuraciones.
      \item Son los ``relojes'' del Juego de la Vida.
    \end{itemize}
  \end{block}

  \vspace{0.3cm}

  \begin{center}
    \textbf{Demo en vivo}

    \demo{beehive} \quad \demo{loaf} \quad \demo{boat}
    \quad \demo{toad} \quad \demo{beacon}
  \end{center}
\end{frame}

% ================================================================
% . NAVES ESPACIALES -- Israel
% ================================================================
\begin{frame}{Patrones que se mueven}

  \vspace{0.4cm}

  \begin{itemize}
    \item Las \textbf{naves espaciales} se desplazan por la grilla
          recuperando su forma original cada cierto número de pasos.
    \item El \textbf{glider} es la nave más pequeña: se mueve en diagonal
          y recorre 1 celda cada 4 generaciones (c/4).
    \item Las naves \textbf{LWSS, MWSS y HWSS} son ortogonales y viajan
          al doble de velocidad (c/2).
  \end{itemize}

  \begin{block}{¿Para qué sirven?}
    Transmiten \textbf{información} a distancia.
    El Gosper Glider Gun las usa como ``mensajeras''.
  \end{block}

  \vspace{0.3cm}

  \begin{center}
    \textbf{Demo en vivo}

    \demo{glider} \quad \demo{lwss} \quad \demo{mwss} \quad \demo{hwss}
  \end{center}
\end{frame}

% ================================================================
% 5. OSCILADORES COMPLEJOS Y CAÑONES -- Jorge
% ================================================================
\begin{frame}{Osciladores grandes y generadores de naves}

  \vspace{0.4cm}

  \begin{itemize}
    \item El \textbf{pulsar} es uno de los osciladores más icónicos:
          período 3, 48 células, simetría perfecta.
    \item El \textbf{pentadecathlon} tiene período 15 y solo 12 células.
    \item El \textbf{Gosper Glider Gun} fue el primer cañón descubierto
          (1970). Produce un glider nuevo cada 30 generaciones.
  \end{itemize}

  \begin{alertblock}{Dato curioso}
    Conway ofreció \$50 a quien demostrara crecimiento infinito.
    Bill Gosper lo ganó con este cañón.
  \end{alertblock}

  \vspace{0.3cm}

  \begin{center}
    \textbf{Demo en vivo}

    \demo{pulsar} \quad \demo{pentadecathlon} \quad \demo{gosper\_glider\_gun}
  \end{center}
\end{frame}

% ================================================================
% 6. MATUSALENES Y PARALELISMO -- Miguel
% ================================================================
\begin{frame}{Matusalenes y el poder del paralelismo}

  \vspace{0.4cm}

  \begin{block}{Matusalenes}
    \begin{itemize}
      \item Patrones pequeños que tardan \textbf{miles de generaciones} en estabilizarse.
      \item \texttt{r\_pentomino}: 5 células $\rightarrow$ 1,103 generaciones $\rightarrow$ 116 células.
      \item \texttt{diehard}: desaparece completamente tras 130 pasos.
      \item \texttt{acorn}: 7 células $\rightarrow$ \textbf{5,206 generaciones}.
    \end{itemize}
  \end{block}

  \begin{block}{¿Por qué paralelismo?}
    \begin{itemize}
      \item La versión \texttt{juego\_vida\_proyecto.py} usa \texttt{multiprocessing}.
      \item Divide un tablero de \textbf{200×200 celdas} entre todos los núcleos del CPU.
      \item Cada proceso calcula su franja de filas de forma independiente.
    \end{itemize}
  \end{block}

  \vspace{0.2cm}

  \begin{center}
    \textbf{Demo en vivo}

    \demo{r\_pentomino} \quad \demo{diehard} \quad \demo{acorn}

    \small También: \texttt{python juego\_vida\_proyecto.py}
  \end{center}
\end{frame}

% ================================================================
% 7. REGLAS ALTERNATIVAS -- Andrés
% ================================================================
\begin{frame}{Más allá de Conway: mismas células, otras reglas}

  \vspace{0.4cm}

  \begin{itemize}
    \item Conway usa la regla \textbf{B3/S23}. Pero hay muchas otras.
    \item Cambiar los números de nacimiento y supervivencia produce
          \textbf{universos completamente distintos}.
  \end{itemize}

  \vspace{0.3cm}

  \begin{columns}[T]
    \column{0.48\textwidth}
    \begin{block}{HighLife (B36/S23)}
      \begin{itemize}
        \item Muy parecido a Conway pero...
        \item ¡Tiene un \textbf{replicador} natural!
        \item Un patrón que se copia a sí mismo.
      \end{itemize}
    \end{block}

    \column{0.48\textwidth}
    \begin{block}{Day \& Night (B3678/S34678)}
      \begin{itemize}
        \item Es \textbf{simétrico}: si inviertes vivo por muerto,
              se comporta exactamente igual.
        \item El día y la noche son indistinguibles.
      \end{itemize}
    \end{block}
  \end{columns}

  \vspace{0.3cm}

  \begin{center}
    \textbf{Demo en vivo}

    \demo{highlife} \quad \demo{seeds} \quad \demo{daynight}

    \small Probar con el mismo patrón: \texttt{gol.load\_pattern("glider")}
    \texttt{gol.set\_rule("highlife")}
  \end{center}
\end{frame}
% ================================================================
% APLICACIONES
% ================================================================
\begin{frame}{Aplicaciones reales}

\begin{itemize}
    \item Simulación de sistemas biológicos.
    \item Modelado de propagación de epidemias.
    \item Redes complejas y tráfico.
    \item Computación distribuida.
    \item Inteligencia artificial y sistemas emergentes.
\end{itemize}

\vspace{0.5cm}

\begin{alertblock}{Idea central}
El Juego de la Vida demuestra cómo reglas locales simples
pueden producir comportamientos globales complejos.
\end{alertblock}

\end{frame}
% ================================================================
% 8. CIERRE -- Todos
% ================================================================
\begin{frame}{Lo que aprendimos}
  \textbf{Equipo Rosa}

  \vspace{0.4cm}

  \begin{itemize}
    \item Con solo \textbf{4 reglas} emerge una complejidad infinita.
    \item El Juego de la Vida es un \textbf{laboratorio} para explorar:
          patrones, movimiento, autorreplicación, caos.
    \item Implementamos \textbf{3 versiones} del mismo algoritmo:
          interactiva, optimizada y paralela.
    \item Cada versión enseña algo distinto sobre cómo estructurar
          y ejecutar código eficientemente.
    \item Cambiar las reglas (\texttt{automata.py}) abre la puerta
          a \textbf{infinitos universos} por explorar.
  \end{itemize}

  \vspace{0.5cm}

  \begin{center}
    \Large \textbf{¿Preguntas?}

    \vspace{0.3cm}
    \normalsize
    Repositorio: \texttt{Equipo-Rosa-PPyC-2026}

    Programación Paralela y Concurrente -- UNAM 2026
  \end{center}
\end{frame}

\end{document}